\section{Des écrans exacts à l'interface physique de la théorie M}
\label{sec:interfaz-teoria-m-rev8}
\index{théorie M!interface physique}
\index{dualité T}
\index{théorie des cordes}
\index{CFT}

La première partie a construit, dans des domaines mathématiques explicites, les
réductions
\[
12\longrightarrow11\longrightarrow10,
\qquad
26\longrightarrow24,
\]
et l'involution de l'impulsion, de l'enroulement et du rayon
\[
(m,w,R)\longleftrightarrow(w,m,R^{-1}).
\]
Ces résultats ne constituent pas encore une théorie physique des cordes ni une action
de supergravité. Ils constituent les écrans et les équivalences qu'une
réalisation dimensionnelle doit respecter.

L'interface est typée au moyen d'une application déclarée
\[
\mathcal R_M:\mathcal X_{\rm HMT}^{\rm enr}
\longrightarrow
\mathcal X_{11}^{\rm phys},
\]
dont le codomaine doit porter, au minimum, le champ métrique \(g_{MN}\), la
trois-forme \(C_{MNP}\), le gravitino \(\psi_M\), son action et les
transformations de supersymétrie. Une spécialisation aux cordes exige
en outre d'identifier le rayon physique, la tension, les secteurs d'impulsion et d'
enroulement et l'espace des états sur lequel se réalise la dualité
T. Une lecture conforme ou CFT requiert sa propre algèbre d'opérateurs,
sa graduation et sa correspondance état--opérateur.

La connexion positive avec la branche exceptionnelle n'est pas une flèche nue
Monstre\(\to\)théorie M. Les deux prolongements conservent des données du même
état enrichi --- phase, graduation, orientation, couture et mémoire ---,
mais emploient des foncteurs distincts. Dans la branche exceptionnelle, ces données atteignent
\(V^\natural\), le Monstre et Moonshine ; dans la branche dimensionnelle, elles alimentent
les écrans et, uniquement par \(\mathcal R_M\), une éventuelle réalisation de
champs et de cordes.

\begin{definition}[Admissibilité interne d'une compactification]
\label{def:admisibilidad-compactificacion-hmt-rev3}
Une compactification candidate est représentée par un triplet
\[
 (x_\infty,\kappa,\mathcal U),
\]
où \(x_\infty\in\mathcal X_{\rm HMT}^{\rm enr}\) est une section
survivante,
\[
 \kappa:\mathcal X_{\rm HMT}^{\rm enr}\dashrightarrow\mathcal M
\]
est l'application, éventuellement partielle, vers un espace de modules déclaré
\(\mathcal M\), et
\(\mathcal U\) transporte la métrique, l'orientation, le registre et les données de frontière.
Nous définissons
\[
 \operatorname{Adm}_{\rm HMT}(x_\infty,\kappa,\mathcal U)=1
\]
si les quatre conditions suivantes sont réunies :
\begin{enumerate}
 \item la section possède un prolongement compatible à toute profondeur ;
 \item les dualités déclarées sont compatibles au point considéré :
       \[
        \kappa(\mathsf T_{\rm HMT}x_\infty)
        =\mathsf T_{\mathcal M}(\kappa(x_\infty)),
       \]
       lorsque les deux membres sont définis ;
 \item le spectre descend au quotient de compactification ;
 \item les observables conservent leurs domaines, leurs codomaines et leurs types sous
       \(\mathcal U\).
\end{enumerate}
\end{definition}

Ce prédicat est un critère interne d'admissibilité et non une
identification automatique au \emph{Swampland} physique. Les critères de
distance, de tour, de gravité faible et d'absence de symétries globales, ainsi que
les conditions géométriques d'une compactification précise, s'évaluent
ensuite sur \(\operatorname{im}\kappa\) et requièrent les applications
métriques, les champs et les échelles correspondants. La définition précédente fixe l'interface qu'
une réalisation doit préserver ; elle ne démontre pas à elle seule l'existence d'
une compactification physique ni une conjecture du Swampland.

\begin{statusbox}[title={Statut de l'interface}]
Les écrans, les quotients et l'involution T-duale sont des résultats
mathématiques exacts. Le dictionnaire complet vers
\((g_{MN},C_{MNP},\psi_M)\), l'action, la supersymétrie et une CFT physique est
une interface typée qui doit être vérifiée dans chaque réalisation. Cette frontière
ne diminue pas le résultat mathématique et n'autorise pas à le présenter comme une dynamique
physique déjà achevée.
\end{statusbox}
